% Requires: \usepackage{amsmath,amssymb}
\subsection{The natural matrix generators}

We use the ordering of the side-pairing matrices in the source data.  If the
original labels are denoted by $M_j$, then our numbering is
\[
\begin{aligned}
&h_{1}=M_{0},\quad h_{2}=M_{1},\quad h_{3}=M_{2},\quad h_{4}=M_{3}, \\
&h_{5}=M_{4},\quad h_{6}=M_{6},\quad h_{7}=M_{8},\quad h_{8}=M_{9}, \\
&h_{9}=M_{10},\quad h_{10}=M_{11},\quad h_{11}=M_{13},\quad h_{12}=M_{14}, \\
&h_{13}=M_{16},\quad h_{14}=M_{20},\quad h_{15}=M_{25},\quad h_{16}=M_{26}.
\end{aligned}
\]
Thus
\[
H_{\mathrm{mat}}=\langle h_1,\ldots,h_{16}\rangle
\leq \operatorname{SO}(5,1;\mathbb Z).
\]

\begin{description}
\item[$h_{1}$] 
\[
h_{1}=\begin{pmatrix}
1 & 0 & -1 & 0 & 0 & 1 \\
0 & 0 & 0 & 1 & 0 & 0 \\
0 & 0 & -1 & 0 & -1 & 1 \\
1 & 0 & 0 & 0 & -1 & 1 \\
0 & 1 & 0 & 0 & 0 & 0 \\
1 & 0 & -1 & 0 & -1 & 2
\end{pmatrix}.
\]
\item[$h_{2}$] 
\[
h_{2}=\begin{pmatrix}
1 & 0 & 0 & 0 & 0 & 0 \\
0 & 1 & 0 & 0 & 0 & 0 \\
0 & 0 & 1 & 0 & 0 & 0 \\
0 & 0 & 0 & 1 & 0 & 0 \\
0 & 0 & 0 & 0 & 1 & 0 \\
0 & 0 & 0 & 0 & 0 & 1
\end{pmatrix}.
\]
\item[$h_{3}$] 
\[
h_{3}=\begin{pmatrix}
0 & 1 & 0 & 0 & 0 & 0 \\
1 & 0 & 0 & 0 & 0 & 0 \\
0 & 0 & 1 & -1 & 0 & 1 \\
0 & 0 & 1 & 0 & -1 & 1 \\
0 & 0 & 0 & -1 & -1 & 1 \\
0 & 0 & 1 & -1 & -1 & 2
\end{pmatrix}.
\]
\item[$h_{4}$] 
\[
h_{4}=\begin{pmatrix}
-1 & 0 & 0 & 0 & -1 & 1 \\
0 & 0 & 1 & 0 & 0 & 0 \\
0 & 1 & 0 & 0 & 0 & 0 \\
0 & 0 & 0 & 1 & -1 & 1 \\
-1 & 0 & 0 & 1 & 0 & 1 \\
-1 & 0 & 0 & 1 & -1 & 2
\end{pmatrix}.
\]
\item[$h_{5}$] 
\[
h_{5}=\begin{pmatrix}
0 & -1 & -1 & 0 & 0 & 1 \\
0 & 0 & 1 & 1 & 0 & -1 \\
1 & 0 & 0 & 0 & 0 & 0 \\
0 & 0 & 0 & 0 & -1 & 0 \\
0 & -1 & 0 & -1 & 0 & 1 \\
0 & -1 & -1 & -1 & 0 & 2
\end{pmatrix}.
\]
\item[$h_{6}$] 
\[
h_{6}=\begin{pmatrix}
0 & -1 & -1 & 0 & 0 & 1 \\
0 & 0 & 1 & 1 & 0 & -1 \\
1 & 0 & 0 & 0 & 0 & 0 \\
0 & 0 & -2 & 0 & -1 & 2 \\
0 & -1 & -2 & -1 & -2 & 3 \\
0 & -1 & -3 & -1 & -2 & 4
\end{pmatrix}.
\]
\item[$h_{7}$] 
\[
h_{7}=\begin{pmatrix}
-1 & -2 & 0 & -1 & -2 & 3 \\
0 & -2 & 0 & 0 & -1 & 2 \\
-1 & -2 & -1 & 0 & -2 & 3 \\
0 & -1 & 0 & 0 & -2 & 2 \\
0 & -2 & -1 & -1 & -2 & 3 \\
-1 & -4 & -1 & -1 & -4 & 6
\end{pmatrix}.
\]
\item[$h_{8}$] 
\[
h_{8}=\begin{pmatrix}
0 & 0 & 0 & 0 & 1 & 0 \\
1 & 0 & 0 & 1 & 0 & -1 \\
0 & 1 & 0 & 1 & 0 & -1 \\
0 & 0 & 1 & 0 & 0 & 0 \\
-1 & -1 & 0 & 0 & 0 & 1 \\
-1 & -1 & 0 & -1 & 0 & 2
\end{pmatrix}.
\]
\item[$h_{9}$] 
\[
h_{9}=\begin{pmatrix}
0 & -1 & 0 & -1 & 0 & 1 \\
0 & 2 & 2 & 1 & 1 & -3 \\
0 & -1 & 0 & 0 & -1 & 1 \\
1 & 0 & 0 & 0 & 0 & 0 \\
0 & -2 & -1 & 0 & 0 & 2 \\
0 & -3 & -2 & -1 & -1 & 4
\end{pmatrix}.
\]
\item[$h_{10}$] 
\[
h_{10}=\begin{pmatrix}
-2 & -1 & 0 & -1 & -2 & 3 \\
0 & 0 & 0 & 1 & 1 & -1 \\
0 & -1 & 0 & 0 & -1 & 1 \\
-1 & 0 & 0 & 0 & -2 & 2 \\
0 & 0 & 1 & 0 & 0 & 0 \\
-2 & -1 & 0 & -1 & -3 & 4
\end{pmatrix}.
\]
\item[$h_{11}$] 
\[
h_{11}=\begin{pmatrix}
0 & 0 & 0 & 0 & 1 & 0 \\
1 & 0 & 0 & 1 & 0 & -1 \\
-2 & -1 & -2 & -1 & 0 & 3 \\
-2 & 0 & -1 & 0 & 0 & 2 \\
-1 & -1 & 0 & 0 & 0 & 1 \\
-3 & -1 & -2 & -1 & 0 & 4
\end{pmatrix}.
\]
\item[$h_{12}$] 
\[
h_{12}=\begin{pmatrix}
-1 & -2 & 0 & 0 & 0 & 2 \\
2 & 1 & 0 & 0 & 0 & -2 \\
0 & 0 & 1 & 0 & 0 & 0 \\
0 & 0 & 0 & 1 & 0 & 0 \\
0 & 0 & 0 & 0 & 1 & 0 \\
-2 & -2 & 0 & 0 & 0 & 3
\end{pmatrix}.
\]
\item[$h_{13}$] 
\[
h_{13}=\begin{pmatrix}
1 & 0 & 0 & -1 & 0 & 1 \\
0 & 0 & 0 & 0 & -1 & 0 \\
1 & 0 & -1 & 0 & 0 & 1 \\
0 & -1 & 0 & 0 & 0 & 0 \\
0 & 0 & -1 & -1 & 0 & 1 \\
1 & 0 & -1 & -1 & 0 & 2
\end{pmatrix}.
\]
\item[$h_{14}$] 
\[
h_{14}=\begin{pmatrix}
0 & 0 & -1 & -1 & 0 & 1 \\
0 & -1 & 1 & 0 & 0 & -1 \\
-1 & 1 & -1 & -1 & 1 & 2 \\
-1 & 0 & -1 & 0 & 0 & 1 \\
0 & 0 & -1 & 0 & 1 & 1 \\
-1 & 1 & -2 & -1 & 1 & 3
\end{pmatrix}.
\]
\item[$h_{15}$] 
\[
h_{15}=\begin{pmatrix}
-1 & 2 & -1 & -2 & 0 & 3 \\
0 & -2 & 0 & 1 & 0 & -2 \\
0 & 2 & -1 & -2 & -1 & 3 \\
-1 & 2 & 0 & -2 & -1 & 3 \\
0 & 1 & 0 & -2 & 0 & 2 \\
-1 & 4 & -1 & -4 & -1 & 6
\end{pmatrix}.
\]
\item[$h_{16}$] 
\[
h_{16}=\begin{pmatrix}
0 & 1 & -2 & 0 & 0 & 2 \\
-1 & 0 & 0 & 0 & 0 & 0 \\
0 & 0 & -1 & -1 & 0 & 1 \\
0 & 0 & -1 & 0 & -1 & 1 \\
0 & 2 & -2 & -1 & -1 & 3 \\
0 & 2 & -3 & -1 & -1 & 4
\end{pmatrix}.
\]
\end{description}

\subsection{The Reidemeister--Schreier generators}

For homomorphism~15, the images of the generators of $T$ are
\[
\phi(x)=h_{15},\qquad \phi(y)=h_{16}.
\]
The two generators in the simplified Reidemeister--Schreier presentation
\[
P=\langle F_1,F_2\mid R_1,\ldots,R_6\rangle
\cong H_{\mathrm{mat}}
\]
are represented by the following words in the natural matrix generators:
\begin{align*}
F_1={}&
h_{15}^{-3}h_{16}^2h_{15}^{-2}h_{16}h_{15}^2
h_{16}^{-1}h_{15}^{-3}h_{16}h_{15},\\
F_2={}&
h_{16}h_{15}^{-2}h_{16}^{-2}h_{15}^{-1}h_{16}^2
h_{15}h_{16}h_{15}^2h_{16}^{-1}h_{15}h_{16}^{-1}.
\end{align*}
Exact multiplication gives
\[
F_1=
\begin{pmatrix}
0&1&0&0&-1&1\\
0&0&1&0&-1&1\\
0&0&0&1&0&0\\
0&1&1&0&0&1\\
1&0&0&0&0&0\\
0&1&1&0&-1&2
\end{pmatrix},
\qquad
F_2=
\begin{pmatrix}
1&0&0&1&0&-1\\
0&0&0&0&1&0\\
-1&0&-1&0&0&1\\
0&1&0&0&0&0\\
0&0&-1&-1&0&1\\
-1&0&-1&-1&0&2
\end{pmatrix}.
\]

\subsection{Verification checks}

All matrix calculations below were performed exactly over $\mathbb Q$ (in
fact, all displayed matrices have integral entries).
\begin{enumerate}
\item For every natural generator, the Lorentz-form identity
\[
h_i^\mathsf{T}Jh_i=J,
\qquad
J=\operatorname{diag}(1,1,1,1,1,-1),
\]
was verified.

\item The seven defining relators $r_k$ of $T$ were evaluated at
$x=h_{15}$ and $y=h_{16}$, and each gave
\[
r_k(h_{15},h_{16})=I_6.
\]
This verifies that homomorphism~15 is well defined.

\item Each of $h_1,\ldots,h_{16}$ was expressed as an explicit word in
$h_{15}^{\pm1}$ and $h_{16}^{\pm1}$, and every resulting matrix
identity was checked exactly.  This verifies that homomorphism~15 is
surjective.

\item Each $h_i$ was expressed as a word in the six standard reflections of
the Coxeter group $\Gamma^5$, and the corresponding matrix identity was
checked exactly.  GAP then verified that the resulting subgroup has index
$3840$ in $\Gamma^5$.

\item Reidemeister--Schreier rewriting, followed by Tietze simplification,
produced
\[
P=\langle F_1,F_2\mid R_1,\ldots,R_6\rangle.
\]
The two displayed words for $F_1,F_2$ were checked by exact matrix
multiplication, and all six relator products were checked to satisfy
\[
R_i(F_1,F_2)=I_6\qquad(1\leq i\leq6).
\]

\item Finally, under the proposed inverse $\psi:P\to T$, all six words
$\psi(R_i)$ reduced to the identity using the defining relators of $T$.
The two composite checks
\[
\psi(\phi(x))=x,
\qquad
\psi(\phi(y))=y
\]
were also verified.  Hence $\psi\circ\phi=\operatorname{id}_T$; together
with surjectivity, this proves $T\cong H_{\mathrm{mat}}$.
\end{enumerate}
